\scanpage{6}
\[
\begin{array}{ll}
V_{\ell_{\bullet}}:\ell_{\bullet}\mathbin{-}\ell_{\bullet}^{(p)}&F_{\ell_{\bullet}}=0:\ell_{\bullet}^{(p)}\mathbin{-}\ell_{\bullet}\\
V_{\ell'_{\bullet}}:\ell'_{\bullet}\mathbin{-}\ell_{\bullet}'{}^{(p)}&F_{\ell'_{\bullet}}=0:\ell_{\bullet}'{}^{(p)}\mathbin{-}\ell'_{\bullet}
\end{array}
\]
1) le triangle $(\mathrm {III})$ est compatible avec $(F,V)$-structure : \uuline{\unsure{autos}} !

2) L'isom $\varphi^*(\mathrm {III})=\mathrm I$ compatible avec $(F-V)$-structures.

\deleted{\unsure{un}}\qquad\deleted{$\unsure{\Delta^{*\,(p)}\longrightarrow\Delta^*}$}

\uline{1)}
\diagram{\begin{tikzcd}
\Delta^{*p}\arrow[rr,"F_{\Delta^*}"]\arrow[dr,"\deleted{\unsure{*}}\pi_0^{(p)}"']&&\Delta^*\\
&\check\ell'{}^{(p)}\arrow[ur,"\varepsilon^*(i_{\mathrm {II}})"']&
\end{tikzcd}}
\[
\begin{gathered}\pi_0^{(p)}\\\Vert\\\varepsilon^*(\varphi^*(\pi_0))\\\Vert\\\varepsilon^*(\pi_{\mathrm I})\end{gathered}
\]
\margin{impl que $i_0$ comp. avec $F$}
\diagram{\begin{tikzcd}
\Delta^*\arrow[rr,dash,"V_{\Delta^*}"]\arrow[dr,"\varepsilon^*(\pi_{\mathrm {II}})"']&&\Delta^{*\,(p)}\\
&\ell_{\bullet}^{\unsure{(p)}}[-1]\arrow[ur,"i_0^{(p)}=\varepsilon^*\varphi^*(i_0)=\varepsilon^*(i_{\mathrm I})"']&
\end{tikzcd}}
\margin{impl que $\pi_0$ compatible avec $V$}

$i_0$ compatible avec $V$:
\diagram{\begin{tikzcd}
\ell_{\bullet}[-1]\arrow[r,"i_0"]&\Delta^*\arrow[d,"V_{\Delta^*}"]\arrow[dr,"\varepsilon^*(\pi_{\mathrm {II}})"]&\\
\ell_{\bullet}[-1]^{(p)}\arrow[r,dash,"i_0^{(p)}"]&\Delta^{*\,(p)}&\ell_{\bullet}^{(p)}[-1]\arrow[l,"i_0^{(p)}"]
\end{tikzcd}}
\[
\varepsilon^*(\pi_{\mathrm {II}})i_0\ne V_{\ell_{\bullet}}[-1]
\]
\margin{OK}

on voit de $=$ que $\pi_0$ compatible avec $F$.

$\partial_0$ compatible avec $F$
\diagram{\begin{tikzcd}
\check\ell'_{\bullet}{}^{(p)}\arrow[r,"\pi_0\varepsilon^*(i_{\mathrm {II}})"]\arrow[d,"\partial_0^{(p)}"']&\check\ell'_{\bullet}\arrow[d,"\partial_0"]\\
\ell_{\bullet}[-1]^{(p)}\arrow[r,"0"]&\ell_{\bullet}[-1]
\end{tikzcd}}
\margin{$\unsure{1}$\qquad $\underbrace{\partial_0\pi_0}\varepsilon^*(i_{\mathrm {II}})=0$}

on voit de $=$ $\partial_0$ compatible avec $V$.

2) A prouver que
\[
\left\{\begin{aligned}
\varphi^*(V_{\ell_{\bullet}}[-1])&=V_{[\mathcal M^{p}\xrightarrow{F}\mathcal M]}\\
\varphi^*(F_{\check\ell'_{\bullet}}[-1])&=F_{[\mathcal M\xrightarrow{V}\mathcal M^p]}
\end{aligned}\right.
\]
\[
\varphi^*(V_{\ell_{\bullet}}[-1])=\varphi^*(\varepsilon^*(\pi_{\mathrm {II}})i_0)
=\varphi^*\varepsilon^*(\pi_{\mathrm {II}})\varphi^*(i_0)
=\big|\pi_{\mathrm {II}}^{(p)}i_{\mathrm I}=V_{[\mathcal M^{\unsure{(p)}}\mathbin{-}\mathcal M]}\big)
\]
et l'autre relation se prouve de $=$ !!
\mspage{5}
